\documentclass[12pt]{article}
\usepackage[utf8]{inputenc}
\usepackage[T1]{fontenc}
\usepackage{hyperref}
\title{Yinglish Language Resources for K-12 Curriculum Development}
\author{Compiled by Claude Sonnet 5.5 (claude-sonnet-5-5) for Andr\'as Kornai}
\date{}
\begin{document}
\maketitle

\section{Introduction}
The code yib (Glottocode ying1248) is no longer a valid ISO 639-3 language. It was retired and merged into English (eng) under Change Request 2006-044, effective 2007-07-18. The request judged Yinglish to be a dialect of English and entirely intelligible with it. Glottolog keeps the entry for bookkeeping only. No separate Yinglish language community or school curriculum turned up in search, so the resources below are administrative records only.

\section{Primary Educational Resources}
None found. English-language curricula apply instead.

\section{Linguistic and Reference Materials}
\begin{itemize}
\item Glottolog bookkeeping entry for Yinglish (yib): \url{https://glottolog.org/resource/languoid/iso/yib}
\item Glottolog languoid page (ying1248): \url{https://glottolog.org/resource/languoid/id/ying1248.bigmap.html}
\item ISO 639:y overview (Wikipedia): \url{https://en.wikipedia.org/wiki/ISO_639:y}
\end{itemize}

\section{Religious and Community Resources}
None found.

\section{Digital and Audio Resources}
None found.

\section{Curriculum Development Context}
The code is retired and the variety is treated as English, so no native-language K-12 curriculum can be built for it. Standard English materials are the relevant resources.
\end{document}